\section*{M13: Unconditional quantitative envelope for the independent-law $k=3$ chain}
\label{sec:m13-unconditional-chain-envelope}

\paragraph{Setup.}
Consider the ordered binary $k=3$ chain with unit-norm leaves, independent
bilinear laws of operator norm at most $M$, orthogonal projectors, and local
closure-residual cap $\rho=\eta M$.  Let
$C_{3,\mathrm{ind}}^{P,\mathrm{chain}}(\eta)$ denote the projected-root
constant normalized by $\rho M^2L_T$ for this topology.

\paragraph{Gram-matrix lemma.}
Let $N$ be a linear map with $\|N\|_{\mathrm{op}}\le M$, let $u,v$ be
orthonormal, and let $Q$ be an orthogonal projector.  Put
\[
  y=Nv,\qquad s=\|Qy\|.
\]
Then
\[
  |\langle Nu,Qy\rangle|
  \le M s\sqrt{1-s^2/M^2}.
\tag{1}
\]
Indeed, with $\widetilde N=N/M$ the operator
$K=\widetilde N^*Q\widetilde N$ satisfies $0\preceq K\preceq I$.  On
$\operatorname{span}\{u,v\}$, write $K_{vv}=s^2/M^2$ and
$K_{uu}=a$.  Positivity of $K$ and $I-K$ gives
\[
 |K_{uv}|^2\le a(s/M)^2,qquad
 |K_{uv}|^2\le(1-a)(1-s^2/M^2).
\]
Maximizing the minimum of the two right-hand sides over $0\le a\le1$
occurs at $a=1-s^2/M^2$ and yields (1).  This argument uses no exact
operator-norm saturation.

\paragraph{Theorem.}
For every $0<\eta\le1$,
\[
 C_{3,\mathrm{ind}}^{P,\mathrm{chain}}(\eta)
 \le W_3(\eta):=
 \begin{cases}
 \sqrt{4-3\eta^2}, & 0<\eta\le\sqrt{2/3},\\[2mm]
 \dfrac{2}{\sqrt3\,\eta}, & \sqrt{2/3}\le\eta\le1.
 \end{cases}
\tag{2}
\]

\paragraph{Proof.}
At the first internal node write $D_1\perp R_1$, put
$q=\|D_1\|/M$ and $p=\|R_1\|/M$.  Then
\[
 0\le q\le\eta,qquad p\le\sqrt{1-q^2}.
\]
At the second node freeze the sibling leaf and write
$N(x)=\mu_2(x,c)$.  For the normalized directions
$u=D_1/\|D_1\|$ and $v=R_1/\|R_1\|$, put
$x=N u/M$, $y=N v/M$, $Q=I-P_2$, and
$s=\|Qy\|\le\eta$.  The exact chain error expansion gives propagated
contributions proportional to $q x$ and $pQy$.  The root map is a contraction
after freezing its sibling leaf, so (1) gives
\[
 \frac{E_T^P}{M^3L_T}
 \le
 \sqrt{q^2+p^2s^2+2qps\sqrt{1-s^2}}.
\tag{3}
\]
Here $\|x\|\le1$ was used in the first term, and the Gram lemma controls the
cross term.  The right side increases with $p$, so use
$p=\sqrt{1-q^2}$.  Set $q=\sin\alpha$, $s=\sin\beta$; then its square is
\[
 f(\alpha,\beta)=\sin^2\alpha+\cos^2\alpha\sin^2\beta
 +2\sin\alpha\cos\alpha\sin\beta\cos\beta.
\]
With $u=\tan\alpha$, $v=\tan\beta$,
\[
 f=\frac{(u+v)^2+u^2v^2}{(1+u^2)(1+v^2)}.
\]
The identity
\[
 4(1+u^2)(1+v^2)-3\big((u+v)^2+u^2v^2\big)
 =u^2v^2+u^2+v^2-6uv+4
 \ge(uv-2)^2
\]
shows $f\le4/3$.  If
$\eta\le\sqrt{2/3}$, then
$u,v\le\tan(\arcsin\eta)\le\sqrt2$ and both partial derivatives are
nonnegative on the square, so the maximum is at $q=s=\eta$.  This gives
$f\le\eta^2(4-3\eta^2)$.  For larger $\eta$, the square contains the
global equality point $u=v=\sqrt2$, giving $f\le4/3$.  Dividing (3) by
$\rho M^2L_T=\eta M^3L_T$ proves (2).  Zero $q$ or $p$ cases follow by
continuity. $\square$

\paragraph{Explicit improvement over M9.}
Let $U_3$ be the unconditional M9 envelope.  Then
\[
 C_{3,\mathrm{ind}}^{P,\mathrm{chain}}(\eta)
 \le U_3(\eta)-\delta_{\mathrm{chain}}(\eta),
\]
where $\delta_{\mathrm{chain}}(\eta)>0$ is the following explicit function:
\[
\delta_{\mathrm{chain}}(\eta)=
\begin{cases}
1+\sqrt{1-\eta^2}-\sqrt{4-3\eta^2},
 &0<\eta\le\eta_*,\\[1mm]
\dfrac{\sqrt{1+\eta^2}}{\eta}-\sqrt{4-3\eta^2},
 &\eta_*<\eta\le\sqrt{2/3},\\[3mm]
\dfrac{\sqrt{1+\eta^2}-2/\sqrt3}{\eta},
 &\sqrt{2/3}<\eta\le1,
\end{cases}
\]
with $\eta_*=\sqrt{(\sqrt5-1)/2}$.  Positivity follows respectively from
$2\sqrt{1-\eta^2}(1-\sqrt{1-\eta^2})>0$, from
$3\eta^4-3\eta^2+1>0$, and from
$1+\eta^2>4/3$ on the stated intervals.

\paragraph{Boundary.}
This is an unconditional chain-topology envelope.  Its sharpness is supplied
by the explicit M14 construction in
\texttt{m14\_exact\_chain\_constant.tex}; it does not provide the same
quantitative envelope for branching.
